\section{Product-belief-first: escribir las creencias sobre el comportamiento en la interfaz, los datos y el postentrenamiento}

La segunda vía parte de una afirmación de producto más fuerte: ¿cómo debe trabajar el sistema junto con las personas? Este tipo de objetivo a menudo no puede deducirse de un único benchmark. El equipo necesita primero construir prototipos de interacción que hagan que los usuarios revelen «dónde se parece a un colaborador y dónde se parece a un autocompletado», y después convertir esas diferencias en ejemplos de entrenamiento, reward y eval. La belief de producto no es una preferencia estética que se impone sobre la evidencia, sino una hypothesis que espera ser operacionalizada y refutada.

\lecturefigure{slide-23-product-belief-vision.jpg}{Vía dos: partir de la product belief y la vision, y después lograr que el modelo lo haga}{00:13:39--00:14:28}

\begin{importantbox}{Cómo escribir una product belief que pueda verificarse}
Ni «el modelo debe ser más inteligente» ni «la experiencia debe ser más natural» pueden verificarse directamente. Una belief operativa debe incluir: el usuario y la tarea objetivo, los comportamientos que el modelo debe adoptar o evitar, el control que conserva el usuario, la ruta de recuperación ante fallos, la evidence de éxito y los límites de costo. Por ejemplo, «ser un collaborator de escritura» implica, como mínimo, que el modelo pueda leer la selección, proponer modificaciones locales, explicar los cambios, conservar la voz del autor y permitir deshacerlos paso a paso.
\end{importantbox}

\subsection{Language as a high-end product}

La ponente considera el language en sí como material de producto de alta gama. El lenguaje no es relleno de la interfaz: el tono, la densidad de información, la estructura, el grado de franqueza, las citas y la expresión de la incertidumbre modifican la manera en que el usuario entiende el mundo. El ejemplo del producto de noticias muestra que un mismo context block puede presentarse de formas distintas en la story page, en el carousel de la página de inicio y en el bottom sheet; los hechos subyacentes se comparten, pero la organización superficial debe adaptarse a cada escenario.

\lecturefigure{slide-24-language-high-end-product.jpg}{Language as a high-end product: la organización de la información es en sí misma comportamiento del producto}{00:14:28--00:14:36}

\lecturefigure{slide-25-context-block-multiple-surfaces.jpg}{Context block compartido: los mismos hechos se mantienen coherentes entre surfaces}{00:14:36--00:14:45}

\lecturefigure{slide-26-adaptive-presentation-surfaces.jpg}{La misma información se presenta de forma adaptativa en el carousel en línea, la página de inicio y el bottom sheet}{00:14:45--00:15:39}

Puede denotarse la representación semántica compartida como $z$, la situación del usuario como $u$ y la surface como $s$; la presentación final es
\[
y_s=g_s(z,u).
\]
$g_s$ se encarga de la extensión, el diseño y la forma de interacción, pero no debe alterar en silencio los hechos contenidos en $z$. En ingeniería, conviene separar en capas el factual core, la presentation policy y el user state; de lo contrario, una edición en una surface puede generar incoherencias en otra, o el modelo puede eliminar matices clave para adaptarse a una interfaz breve.

\begin{knowledgebox}{Primer uso: Single source of truth}
Single source of truth no exige que todas las interfaces muestren la misma frase, sino que varias surfaces hagan referencia al mismo objeto de hechos rastreable. El modelo puede resumir, reordenar y reformular, pero cada conclusión debe poder remitirse a la fuente compartida; cuando el usuario actualiza un hecho en un lugar, los demás lugares deben sincronizarse explícitamente o indicar la versión.
\end{knowledgebox}

\subsection{Del Writing IDE a la micro-personalization}

El ejemplo del Writing IDE inserta el modelo de lenguaje en el flujo de trabajo del autor, en lugar de pedirle que copie todo el texto en una ventana de chat. La verdadera dificultad no es generar una frase bonita, sino entender la selección, el contexto, la intención del autor y el alcance de los cambios. La micro-personalization va más allá y exige que el sistema recuerde «cómo trabaja este usuario», pero la personalización no debe convertirse en una manipulación invisible.

\lecturefigure{slide-27-writing-ide-gpt3.jpg}{Writing IDE: el modelo ofrece colaboración local dentro del contexto del documento}{00:15:39--00:15:50}

\lecturefigure{slide-28-language-compressed-knowledge.jpg}{El lenguaje es un artifact altamente comprimido del conocimiento humano y también material para la innovación de interfaces}{00:15:50--00:16:17}

\lecturefigure{slide-29-claude-micro-personalization.jpg}{Micro-personalization en Claude.ai: organizar los accesos posteriores según los temas del historial}{00:16:17--00:16:20}

Si el estado de personalización es $m_u$, una actualización puede escribirse como
\[
m_u' = U(m_u,e_t,\gamma),
\]
donde $e_t$ es la evidence de la interacción actual y $\gamma$ es la intensidad de retención. El producto debe responder: qué eventos pueden registrarse, cuánto tiempo se conservan, cómo puede el usuario consultarlos o eliminarlos y cómo se revierte un recuerdo erróneo. Si $U$ es completamente invisible, el modelo puede consolidar una preferencia ocasional como un perfil de largo plazo.

\begin{warningbox}{Tres errores de la personalización}
Primero, tomar los clics frecuentes como preferencias estables; segundo, trasladar información sensible entre tareas sin que el usuario lo sepa; tercero, reducir la diversidad de puntos de vista con tal de resultar «atento». La micro-personalization debe medir a la vez relevance, surprise, user control y privacy, y no solo la tasa de clics.
\end{warningbox}

\teachervoice{00:13:39--00:16:18, Karina explica, a partir del producto de noticias, el Writing IDE y el carácter comprimido del lenguaje, cómo la product belief antecede al comportamiento del modelo. El docente enfatiza: la verdadera interface innovation suele requerir entrenar al modelo para que entienda nuevos objetos de interacción, y no solo dibujar una UI alrededor de un modelo existente.}

\subsection{Virtual teammate: asincronía, varias personas y acceso a herramientas}

«AI teammate» se usa con frecuencia de forma abusiva como término de marketing. El prototipo de Claude in Slack que presenta la ponente tiene límites de comportamiento más concretos: conversación asíncrona, colaboración entre varias personas, tareas automáticas, resúmenes periódicos del canal y acceso a herramientas. Todavía no es un colega completo, pero al menos reformula teammate, de un adjetivo de personalidad, a un protocolo de trabajo observable.

\lecturefigure{slide-30-claude-slack-product-spec.jpg}{Claude in Slack: generar product spec, objetivos y OKR a partir del historial del canal}{00:16:20--00:17:50}

\begin{knowledgebox}{Términos clave: Assistant, Tool, Agent, Teammate}
\begin{tabularx}{\textwidth}{L{0.17\textwidth} X X}
\toprule
Rol & Mecanismo central & Pregunta principal de verificación \\
\midrule
Assistant & Genera respuestas a solicitudes & Si la ayuda en un solo turno es correcta y clara \\
Tool & Ejecuta operaciones acotadas & Si la entrada y la salida son deterministas y pueden auditarse \\
Agent & Mantiene estado e invoca herramientas & Permisos, plan, recuperación, condiciones de paro \\
Teammate & Coordina de forma continua en flujos con varias personas & Límites de responsabilidad, contexto compartido, traspaso y juicio social \\
\bottomrule
\end{tabularx}
Un mismo modelo puede desempeñar roles distintos; el rol lo definen el entorno y los permisos, no el tono del chat.
\end{knowledgebox}

Un teammate en un canal con varias personas debe mantener, como mínimo, un estado compartido $s_t$, un estado visible para cada persona $p_i$ y un conjunto de permisos $\mathcal{A}_i$. Las acciones ejecutables deben cumplir
\[
a_t\in \mathcal{A}(s_t,p_i,\text{policy}),
\]
en lugar de ejecutarse solo porque el modelo «considera que ayudan». Resumir el canal, escribir una spec y modificar un issue tienen niveles de riesgo distintos y deben tener umbrales de confirmación distintos.

\subsection{Entrenar al modelo para que sea un collaborator}

Canvas no queda terminado con solo conectar un editor de texto al modelo. El trabajo de entrenamiento que repasa la ponente incluye datos objetivo, uso del context, formato de documentos, editing suggestion, rewriting, formatting y output quality. Los resultados en barras de la figura indican que el equipo descompuso la «colaboración» en varias tareas repetibles y después iteró sobre los datos y el postentrenamiento; no demuestran que una puntuación global represente todos los escenarios de escritura.

\lecturefigure{slide-31-human-ai-flexible-affordances.jpg}{Human--AI collaborative affordance: la interfaz debe ampliarse según la tarea}{00:17:50--00:19:36}

\lecturefigure{slide-32-training-collaborator.jpg}{Entrenar al modelo para que sea un collaborator: varios tipos de comportamiento de edición y comparación de calidad}{00:18:21--00:19:41}

\begin{importantbox}{Contract mínimo de comportamiento de un collaborator}
\begin{enumerate}
  \item Leer el artifact actual, la selección, el objetivo del autor y el alcance de las partes que no deben modificarse;
  \item Proponer primero un plan local y después generar un diff comparable;
  \item Conservar las zonas no autorizadas y la voice del autor;
  \item Dar la fuente o la incertidumbre de los cambios que afectan a los hechos;
  \item Permitir aceptar, rechazar, modificar parcialmente y hacer un rollback completo;
  \item Convertir la retroalimentación del usuario en una eval reproducible para la siguiente iteración, en lugar de limitarse a la adaptación en línea.
\end{enumerate}
\end{importantbox}

\begin{lstlisting}[caption={Ciclo de colaboración en documentos con control del usuario}]
state = load(document, selection, user_goal, permissions)
proposal = model.plan_and_edit(state)
diff = isolate_changes(proposal, protected_regions=state.locked)
evidence = verify_facts_and_style(diff, state.sources)
decision = user.review(diff, evidence)
if decision.accepted:
    commit(diff)
else:
    record_failure(state, diff, decision.reason)
\end{lstlisting}

\teachervoice{00:17:50--00:20:25, en clase se preguntó «qué significa exactamente que el modelo sea un collaborator», y se señaló que Canvas requiere entrenamiento específico en context use, formatting, rewriting y editing suggestion. Experiencia práctica: detrás de una buena interfaz de colaboración hay datos de comportamiento y eval, no un cuadro de entrada más grande.}

\subsection{Modular tool composition: Tasks no es un solo prompt}

La pantalla de ChatGPT Tasks combina la activación programada, la generación de un texto, su edición posterior y su guardado. Aquí el sistema no solo genera texto: también necesita un scheduler, un artifact store, permisos y recuperación ante fallos. Conectar varios módulos produce nuevos state mismatch: una tarea puede ejecutarse sobre un contexto desactualizado, el resultado generado puede sobrescribir las modificaciones del usuario y la ejecución repetida puede provocar conflictos.

\lecturefigure{slide-33-chatgpt-tasks-composition.jpg}{ChatGPT Tasks: combinación modular de tareas programadas, el artifact de Canvas y la edición posterior}{00:19:41--00:20:25}

\begin{knowledgebox}{Primer uso: Idempotency y rollback}
\term{Idempotency (idempotencia)} significa que ejecutar repetidamente la misma acción no produce efectos secundarios adicionales; las tareas programadas, en particular, necesitan una idempotency key para evitar que los reintentos de red publiquen o modifiquen archivos más de una vez. \term{Rollback (reversión)} exige que el sistema pueda volver a una versión conocida. Ambos son fundamentos de los productos de agent/tool y no deben añadirse hasta después de que el modelo falle.
\end{knowledgebox}

\subsection{Resumen de la sección}

El núcleo de Product-belief-first es descomponer «cómo queremos que colabore el modelo» en objetos de interfaz, permisos, datos y eval. Los productos de lenguaje, el context compartido, la personalización, el teammate en Slack, el collaborator de Canvas y Tasks muestran lo mismo: cuanto más cerca se está del trabajo real, más se necesitan versiones, permisos, diff, rollback y control del usuario. La siguiente sección aborda el caso más sólido de esta clase: cómo convertir un model behavior difuso en una eval confiable.
